\scanpage{308}
\mspage{7}
quatre avatars possibles
\[
 K^*\qquad K_*\qquad K^{.}\qquad K_.
\]
liés par les relations ci dessus
\[
 \left\{\begin{aligned}
 K^*&=\mathrm{DP}^*(K^{.}) & K^{.}&=\mathrm{ND}^{.}(K^*)\\
 K_*&=\mathrm{DP}_*(K_.) & K_.&=\mathrm{ND}_ .(K_*)
 \end{aligned}\right.
\]
plus les relations
\[
 \left\{\begin{aligned}
 K^*&=(K_*)^\vee, & K_*&=(K^*)^\vee\\
 K^{.}&=(K_.)^\vee, & K_.&=(K^{.})^\vee .
 \end{aligned}\right.
\]
\deleted{\unsure{(et notre remarque tient}}

Ces \deleted{$\unsure{K^*=|K}$} \unsure{notationss|notations} et remarques (i.e. équivalences de catégories) s'étendent aux objets \deleted{\unsure{dépendant}} \inserted{dépendant de} plusieurs indices, p.ex \deleted{$K^*$} une même lettre $K$ (désignant \inserted{\unsure{initialemt|p.ex. initialement}} un objet $K^*{}_*$ de $(A_*)^*$ \inserted{disons)}, désigne également, \deleted{\unsure{au total juste}} quatre avatars distincts
\[
 \left\{\begin{aligned}
 K^*{}_*&\in(A_*)^*, & (K^{.}{}_*)&\in(A_*^{\deleted{\unsure{*}}})^{.}\\
 K^{.}{}_.&\in(A_.)^{.}, &(K^*{}_.)&\in(A_.)^*
 \end{aligned}\right.
\]
avec des notations
\[
 \begin{aligned}
 K^{[n]}{}_{[m]},&\qquad K^n{}_{[m]}\\
 K^n{}_{m},&\qquad K^{[n]}{}_{m}
 \end{aligned}
\]
pour les composantes. De $=$, on aura, si on part p.ex. de $K^{**}\in(A^*)^*$
\[
 \left\{\begin{aligned}
 K^{**}&\in(A^*)^*, &K^{\unsure{.}*}&\in(A^{.})^*\\
 K^{..}&\in(A^{.})^{.}, &K^{*.}&\in(A^*)^{.}
 \end{aligned}\right.
\]
\[
 \left\{\begin{aligned}
 K^{[n][m]},&\qquad K^{n[m]}\\
 K^{nm},&\qquad K^{[n]m}
 \end{aligned}\right.
\]
On notera qu'un objet de $(A_*)^{.}$ est un complexe de cochaines \inserted{$n\longmapsto(K_*{}^n,d^n)$} dans la catégorie $A_*$ des objets $1/2$ simpliciaux de $A$,
% archive sheet number: 309
